\section{Introduction}

The rapid progress of artificial intelligence (AI) and large language models (LLMs) ~\citep{liu2024deepseek, qwen2025technical} has opened new opportunities for developing autonomous agents capable of tool use, reasoning, and task execution in complex environments. Such agents promise to significantly enhance human productivity by automatically interacting with digital systems and real-world platforms~\citep{DBLP:conf/iclr/YaoZYDSN023, DBLP:conf/acl/DingTPWGDSC25,DBLP:journals/fcsc/QuDWCWYXW25}. 
Broadly, the development of agent systems follows two paradigms. One relies on workflow-based orchestration, where developers manually design action sequences or tool-calling rules~\citep{DBLP:conf/icml/DuW024, DBLP:conf/acl/0001YTSLD025}. The other is learning-based optimization, where agents adapt their policies from feedback and interaction~\citep{DBLP:journals/corr/abs-2503-23383,DBLP:journals/corr/abs-2509-13761}. Among learning-based approaches, reinforcement learning (RL) has emerged as a central direction~\citep{DBLP:journals/corr/abs-2509-12867}, offering a principled way to endow agents with adaptive, generalizable behaviors.

Despite this promise, RL-driven agent development remains challenging in practice, especially when adapting to novel environments~\citep{wang2025modeling,DBLP:journals/corr/abs-2509-02547}. 
First, the construction of training tasks tailored for RL is prohibitively expensive. In novel environments, tool functionalities are often unknown, and manually creating tasks with sufficiently diverse trajectories is labor-intensive. 
Second, existing RL methods for LLM-based agents suffer from severe sample inefficiency. Although the RL community has proposed various exploration strategies such as intrinsic rewards~\citep{DBLP:journals/corr/abs-2502-01456} and UCB-based bandit methods~\citep{song2025outcome}, these techniques are often ineffective for long-horizon, tool-augmented agents, where each rollout is computationally and financially expensive. As a result, current LLM agent training still primarily adopts PPO- or GRPO-style policy optimization~\citep{deepseek_math}, whose reliance on massive trajectory sampling leads to approximate brute-force exploration with many redundant rollouts and limited learning value. Consequently, the training of LLM agents remains a laborious and costly paradigm.

Motivated by the rapidly growing semantic understanding and reasoning capabilities of LLMs, a natural question arises: \textbf{why not entrust the model itself with greater autonomy in driving its own learning process?} Instead of relying on rigid, human-engineered pipelines, we envision an agent system where the LLM actively guides exploration, task generation, and performance refinement. To this end, we propose \textbf{\AgentEvolver}, a self-evolving agent system designed to achieve autonomous and efficient capability evolution through environmental interaction. By harnessing the strengths of LLMs, \AgentEvolver establishes a self-training loop that continuously improves agent competence through direct environment interaction.


Specifically, \AgentEvolver aligns with the typical training flow—i) \emph{from environments to tasks}, ii) \emph{from tasks to trajectories}, and iii) \emph{from trajectories to agent policy}. 
As illustrated in Figure~\ref {fig:agentevolver_overview}, this self-evolution process comprises three synergistic mechanisms: \texttt{self-questioning}, \texttt{self-navigating}, and \texttt{self-attributing}:
\begin{itemize}
    \item \textbf{\texttt{Self-questioning}:} enables curiosity-driven exploration, allowing the LLM to autonomously generate tasks by probing the environment's state-action space and discovering functional boundaries. This mechanism reduces dependence on costly handcrafted datasets.
    \item \textbf{\texttt{Self-navigating}:} improves exploration efficiency by reusing and generalizing from past experiences. Through hybrid policy learning and trajectory guidance, the agent achieves more targeted and efficient task completion.
    \item \textbf{\texttt{Self-attributing}:} enhances sample efficiency by allocating fine-grained rewards along long trajectories. Instead of uniformly attributing outcomes as in typical GRPO methods, the LLM infers the respective contribution of intermediate states and actions, enabling more fine-grained and effective learning.
\end{itemize}
Together, these mechanisms constitute an integrated framework for agent self-evolution, systematically addressing the challenges of task scarcity, inefficient exploration, and low sample utilization. By shifting the training initiative from human-engineered pipelines to LLM-guided self-improvement, \AgentEvolver establishes a new paradigm that paves the way toward scalable, cost-effective, and continually improving intelligent systems.


Beyond the conceptual framework, we further introduce a practical infrastructure with a set of developer-oriented features. In particular, it ensures \emph{environmental compatibility} through standardized interfaces, enabling seamless interaction with diverse external environments. A unified \emph{Context Manager} then controls the agent’s multi-turn interaction logic under a consistent interface. It also integrates with reinforcement learning infrastructures such as \texttt{veRL} \citep{sheng2024hybridflow}, supporting efficient policy optimization and parameter updates. Moreover, the modular architecture facilitates secondary development: individual components can be optimized, extended, or redesigned to meet evolving research and application demands. In this way, \AgentEvolver serves both as a methodological contribution and as a practical foundation for building the next generation of self-evolving agent systems.

In the following sections, we first introduce the problem formulation of this work in Section~\ref{sec:formulation}.
We then present the three core mechanisms of the proposed framework:
Section~\ref{sec:self-questioning} presents the \texttt{self-questioning} module for curiosity-driven task generation; 
Section~\ref{sec:self-navigating} details the \texttt{self-navigating} strategy for efficient exploration; 
and Section~\ref{sec:self-attribution} explains the \texttt{self-attributing} mechanism for fine-grained credit assignment. 
Section~\ref{sec:infrastructure} then describes the supporting infrastructure. 
We evaluate the effectiveness of \AgentEvolver in Section~\ref{sec:experiments}, and finally conclude the report in Section~\ref{sec:conclusion} by summarizing the key findings and discussing future research directions.

\begin{figure}
  \centering
  \includegraphics[width=1.0\textwidth]{figs/flowchart1103.pdf}
  \caption{Overview of the \AgentEvolver framework. The self-evolving process is driven by three synergistic mechanisms: \textbf{\texttt{Self-questioning}} for autonomous task generation, \textbf{\texttt{Self-navigating}} for experience-guided exploration, and \textbf{\texttt{Self-attributing}} for fine-grained credit assignment.}
  \label{fig:agentevolver_overview}
\end{figure}